\section{Front wheeler on MLPs}
\label{sec:exp2}

\paragraph{Conditions.} MLPs with four hidden layers of width 512, trained
from scratch with SGD (momentum $0.9$, learning rate $0.05$, gradient-norm
clipping $1.0$) on MNIST (20 epochs) and CIFAR-10 (30 epochs); batch size
128; 3 seeds. \fw{} uses $C{=}50$, $B{=}50$, $\kappa_{\max}{=}2000$
throughout, with $\tau$ swept over $\{0.005, 0.01, 0.02, 0.05\}$ (MNIST) /
$\{0.005, 0.02, 0.05\}$ (CIFAR-10). Baselines and their parameter sweeps
follow the protocol of Section~\ref{sec:framework}; the complete population
of results at every budget is Table~\ref{tab:exp2}
(Appendix~\ref{app:tables}).

\paragraph{Adaptive bursts at a matched low budget.} The claim under test is that \emph{when} deep steps are
spent matters at scarce budgets. At the ${\sim}4$--$5\%$ deep-step tier,
the strongest fixed schedules reach $95.7$ (MNIST, \lpft{}($0.95$)) and
$47.1$ (CIFAR-10); \fw{} reaches $97.1$ and $51.6$ at its best threshold,
and the comparison does not hinge on threshold selection---the \emph{worst}
$\tau$ still leads the strongest matched fixed baseline by $+0.9$ and
$+3.6$. The margins grow with task difficulty
(Figures~\ref{fig:pareto}--\ref{fig:pareto-cifar}). The answer reverses at
generous budgets: with $50\%$ deep steps, one-switch \lpft{}($0.5$) is the
best schedule on both datasets ($98.0$ / $54.2$), consistent with its
design regime \citep{kumar2022finetuning}---adaptivity pays precisely when
the budget is scarce. (\lpft{} was designed for pretrained trunks; from
scratch it is an imperfect fit that we use as the standard one-switch
baseline.)

\begin{figure}[t]
\centering
\includegraphics[width=0.82\textwidth]{exp2_pareto_mnist}
\caption{MNIST: accuracy vs.\ full-backprop budget (log scale; mean
$\pm$ std, 3 seeds). The \fw{} diamond pools all tested thresholds and
seeds---whiskers span min to max. Gray lines anchor the always-full and
never-full extremes.}
\label{fig:pareto}
\end{figure}

\begin{figure}[t]
\centering
\includegraphics[width=0.82\textwidth]{exp2_pareto_cifar}
\caption{CIFAR-10: as Figure~\ref{fig:pareto}. \textsc{periodic}-$N$ and
one-switch \lpft{} trace nearly identical Pareto curves; the adaptive
schedules sit above both in the scarce-budget regime, by margins that grow
from MNIST to CIFAR-10.}
\label{fig:pareto-cifar}
\end{figure}

\paragraph{Why interleaved fixed schedules fail.} Without
gradient clipping, periodic interleaving collapsed to chance within the
first epoch on all three seeds at $N \in \{20, 50\}$, and on one of three
seeds at $N{=}10$: the front-trained, increasingly confident head
back-propagates error into a trunk it has drifted away from, with trunk
gradient norms growing to ${\sim}20\times$ their initial scale before the
collapse, driving hidden units dead. \fw{}'s \emph{consecutive} bursts
partly mitigate this (the trunk re-equilibrates within a burst), and global
norm-clipping removes it entirely; all reported runs clip at $1.0$. This is
a from-scratch analogue of the feature-distortion argument for \lpft{}
\citep{kumar2022finetuning}, here observable as a training-stability
failure rather than an out-of-distribution effect. Clipping closes the
issue under SGD; Section~\ref{sec:exp3} shows that under Adam it does not,
which is where the schedule's interaction with the optimizer becomes the
story.

\paragraph{The block choice is not the active ingredient.} \bcd{}
shares \fw{}'s closed-loop trigger but spends bursts on the complementary
block (trunk only, head frozen). It ties \fw{} on MNIST ($97.2$ vs.\
$97.1$) and trails by $1.1$ points on CIFAR-10. The scheduling mechanism
therefore transfers across block structures nearly intact; joint bursts add
a modest gain only on the harder task. Note that \bcd{} bursts cannot be
made cheap (Section~\ref{sec:framework}): only the nested scheme converts
its cheap phase into compute savings.

\paragraph{Against progressive freezing.}
The \textsc{ratchet} baseline is compared in backward compute, since its
per-step cost varies as layers freeze (its deep-step column is not
comparable; see the protocol). At matched compute, \textsc{ratchet}($0.1$)
exceeds \fw{} on MNIST consistently but not significantly ($97.3$, its own
accuracy, vs.\ \fw{}'s $97.1$; positive in all three seed pairings, $+0.22$
on average, within one seed-std), \textsc{ratchet}($0.3$) reaches $98.2$ at
a third of full BP's compute, and \fw{} leads on CIFAR-10 ($51.6$ vs.\
$50.8$ at matched compute). We therefore do not claim \fw{} beats all fixed
schedules on easy tasks; the decisive contrast appears on the transformer
(Section~\ref{sec:exp3}), where the same \textsc{ratchet}($0.1$) setting is
catastrophic---frozen-forever layers cannot recover from premature
freezing, while \fw{}'s bursts are revocable by construction.

\paragraph{Cost measurement.} We report budgets in backward
multiply--accumulates (Section~\ref{sec:framework}) because end-to-end wall
clock at this scale is dominated by the harness (data loading, kernel
launch, per-epoch evaluation): all schedules complete in ${\sim}16$s
(MNIST) / ${\sim}20$s (CIFAR-10) per run on an RTX PRO 6000. Isolated
per-step timing confirms the proxy translates once GEMMs dominate: a front
step is $1.4\times$ faster than a full step at width 512/depth 4 (batch
128), rising to $2.5$--$2.7\times$ at widths 4096--8192 and depths 8--16
(batches 512--1024), approaching the ${\sim}3\times$ limit implied by a
backward pass costing ${\sim}2\times$ the forward; the measured ratios
slightly exceed the GEMM-only prediction because the trunk's no-grad
forward also skips activation bookkeeping. Since \fw{} takes ${\sim}96\%$
front steps here, its amortized step cost approaches the front-step cost at
scale.
